The DESI Strong Lens Secondary Target program \citep[among a number of ``secondary target'' programs, see][]{desi2023a} has \ed{so far} observed a \ed{majority} of our lens candidates (Paper~II).
% However, some of the candidates have source redshifts difficult to measure for DESI when the [O II] doublet feature is beyond the optical range, or $z_s > 1.66$.
The success rate of obtaining redshifts for lensing galaxies by DESI is \ed{very} high \ed{($> 90\%$)\footnote{\ed{The reasons for redshift failures are: low signal-to-noise ratio, 
contamination by the presence of source spectral features (typically strong emission lines), spectral reduction issues, and Redrock (the spectral fitting pipeline for DESI) fitting failure.}}}.
These are nearly always bright elliptical galaxies with a number of identifiable features in the optical range,
typically absorption lines
% (e.g. Ca H\&K, G, Mg I, and Na absorption lines) 
and the 4000 \ang break. 
For lensed sources, 
some of the spectra have no clear features,  
likely due to the fact that they are too faint and/or their redshifts place key spectral features outside the optical range. 
Lensed sources are typically star forming galaxies for which the [\ion{O}{2}] doublet emission feature is often used to anchor redshift fits. 
This makes redshift measurements in the optical range for $z_s \gtrsim 1.6$ challenging.\footnote{It is true that for $z_s \gtrsim 2.0$, 
Ly-$\alpha$ will 
be in the optical range, 
though
% In our experience, 
this feature is not always present.}
Thus for some systems obtaining $z_s$
may require NIR spectroscopic follow-up observations. 
Our on-going Keck NIRES program has determined the source redshifts for six of them (Paper~III),
for which $z_s \gtrsim 1.6$.
One of these systems is DESI-165.4754-06.0423. 
Below we present the lens model for this system.